\documentclass[../../../include/open-logic-section]{subfiles}

\begin{document}

\olfileid{sth}{spine}{rank}
\olsection{الرتبة}

لما عرّفنا المراحل بمجموعات~$V_\alpha$، وثبت أن كل مجموعة جزء من
مرحلة ما، أمكن أن نعرّف \emph{رتبة} المجموعة. ورتبة~$A$ في الحدس
هي أول مرحلة تتكون فيها~$A$؛ وأما حدها الدقيق فهو الآتي.

\begin{defn}\ollabel{defnsetrank}
لكل مجموعة~$A$، تكون~$\setrank{A}$ أصغر عدد ترتيبي~$\alpha$ يحقق
$A \subseteq V_\alpha$.
\end{defn}
\begin{prop}\ollabel{ranksexist}
	توجد~$\setrank{A}$ لكل~$A$.
\end{prop}
\begin{proof}
	متروك تمرينًا.
\end{proof}
\begin{prob}
	أثبت~\olref[sth][spine][rank]{ranksexist}.
\end{prob}
\noindent
وبفضل الترتيب الجيد للرتب نستخرج نتائج نافعة، أولها:
\begin{prop}\ollabel{valphalowerrank}
لكل عدد ترتيبي~$\alpha$،
$V_\alpha = \Setabs{x}{\setrank{x} \in \alpha}$.
\end{prop}

\begin{proof}
إذا كان~$\setrank{x} \in \alpha$، كان
$x \subseteq V_{\setrank{x}} \in V_\alpha$؛ فيلزم
$x \in V_\alpha$، لأن~$V_\alpha$ مقتدرة، باستعمال
\olref[Valphabasic]{Valphabasicprops} مرات. وفي العكس، إذا كان
$x \in V_\alpha$، كان~$x \subseteq V_\alpha$، ومن ثم
$\setrank{x} \leq \alpha$. ويبيّن استقراء بسيط متناهٍ في التجاوز أن
$x \notin V_{\setrank{x}}$؛ ويلزم لذلك~$\setrank{x} \neq \alpha$،
ومن ثم~$\setrank{x} \in \alpha$.
\end{proof}
\begin{prob}
	أكمل الاستقراء البسيط المتناهي التجاوز المذكور في
	\olref[sth][spine][rank]{valphalowerrank}.
\end{prob}

\begin{prop}\ollabel{rankmemberslower}
إذا كان~$B \in A$، فإن~$\setrank{B} \in \setrank{A}$.
\end{prop}

\begin{proof}
$A \subseteq V_{\setrank{A}} =
\Setabs{x}{\setrank{x} \in \setrank{A}}$ بموجب
\olref{valphalowerrank}.
\end{proof}
\noindent
ومن هذه الحقيقة نحصل على ضرب من الاستقراء يثبت الأحكام في
\emph{جميع المجموعات}:

\begin{thm}[مخطط الاستقراء بـ$\in$]
لكل صيغة~$\phi$:
\[
	\forall A((\forall x \in A)\phi(x) \lif \phi(A)) \lif \forall A \phi(A).
\]
\end{thm}

\begin{proof}
نبرهن بالمقابل. افترض~$\lnot \forall A \phi(A)$. يقتضي الاستقراء
المتناهي التجاوز
(\olref[ordinals][basic]{ordinductionschema}) وجود مجموعة لا تحقق
$\phi$ ورتبتها أصغر رتبة ممكنة؛ أي يوجد~$A$ بحيث~$\lnot \phi(A)$ و
$\forall x(\setrank{x} \in \setrank{A} \lif \phi(x))$. فإذا كان
$x \in A$، كان~$\setrank{x} \in \setrank{A}$ من
\olref{rankmemberslower}، ومنه~$\phi(x)$. وهكذا
$(\forall x \in A)\phi(x) \land \lnot \phi(A)$.
\end{proof}\noindent
ومعنى ذلك بغير الرموز أن نقول إن~$\phi$ \emph{وراثية} إذا وفقط إذا
كان تحقق~$\phi$ في كل !!{element} من مجموعة موجبًا لتحقق~$\phi$ في
المجموعة نفسها. فمخطط الاستقراء بـ$\in$ يقول: إذا كانت~$\phi$ وراثية،
فإن كل مجموعة تحقق~$\phi$.

ونختم، إلى حين، بحث الرتب بإثبات دعاوى تقدمت الإشارة إليها.

\begin{prop}\ollabel{ranksupstrict}
$\setrank{A} = \supstrict_{x \in A}\setrank{x}$.
\end{prop}

\begin{proof}
ضع~$\alpha = \supstrict_{x \in A}\setrank{x}$. من
\olref{rankmemberslower} ينتج~$\alpha \leq \setrank{A}$. ومن جهة
أخرى، إذا كان~$x \in A$، كان~$\setrank{x} \in \alpha$، ومن ثم
$x \in V_\alpha$ بموجب~\olref{valphalowerrank}. فيكون
$A \subseteq V_\alpha$، أي~$\setrank{A} \leq \alpha$؛ فتتعين
$\setrank{A} = \alpha$.
\end{proof}

\begin{cor}\ollabel{ordsetrankalpha}
لكل عدد ترتيبي~$\alpha$، $\setrank{\alpha} = \alpha$.
\end{cor}

\begin{proof}
افترض، في الاستقراء المتناهي التجاوز، أن
$\setrank{\beta} = \beta$ لكل~$\beta \in \alpha$. عندئذ
$\setrank{\alpha} = \supstrict_{\beta \in \alpha}\setrank{\beta}
= \supstrict_{\beta \in \alpha}\beta = \alpha$ بموجب
\olref{ranksupstrict}.
%	First note that $\setrank{\alpha} \neq \beta$ for any $\beta \in
%	\alpha$. For otherwise we would have $\beta = \setrank{\beta} \in
%	\setrank{\alpha} = \beta$ by \olref{rankmemberslower}, i.e.,
%	$\beta \in \beta$, a contradiction. 
%
%	Now note that $\alpha \subseteq V_\alpha$. For $\alpha =
%	\Setabs{\beta}{\beta \in \alpha}$ by
%	\olref[sfr][ordinals][basic]{ordissetofsmallerord}. And if
%	$\beta \in \alpha$ then $\beta \subseteq V_{\setrank{\beta}} =
%	V_\beta \in V_\alpha$, hence $\beta \in V_\alpha$, by
%	\olref[sfr][spine][Valphabasic]{Valphabasicprops} twice. 
\end{proof}

وأخيرًا نثبت سريعًا ما وُعد به في آخر~\olref[foundation]{sec}، وهو أن
$\ZFminus$ تثبت
$\emph{الانتظام} \Rightarrow \emph{التأسيس}$. ومفهوم «الرتبة»
و\olref{rankmemberslower} مباحان في هذا البرهان؛ لأنهما، كما ذكر في
صدر هذا القسم، يُعرضان باستعمال~$\ZFminus + \text{الانتظام}$.

\begin{prop}[العمل في~$\ZFminus + \text{الانتظام}$]\ollabel{zfminusregularityfoundation}
يصدق التأسيس.
\end{prop}

\begin{proof}
ثبّت~$A \neq \emptyset$، واختر~$B \in A$ ذا أصغر رتبة ممكنة. إذا
كان~$c \in B$، كان~$\setrank{c} \in \setrank{B}$ بموجب
\olref{rankmemberslower}؛ ولذلك~$c \notin A$، باختيار~$B$.
\end{proof}

\end{document}
